% !TeX root = ../papers/nc-p3.tex
\section{標数零の全通常階数15窓の七次までの中心上の自由性}
\label{sec:rank-fifteen-truncated-central-freeness}
\begin{theorem}
標数零の通常な階数15窓のすべてで、四始点に共通な四つの四次中心 $z_1,\ldots,z_4$ を選べる。
各始点 $p$ の長さ $0\le n\le7$ の経路空間 $A_{p,n}$ は次元 $h_n=\binom{n+3}{3}$ を持つ。
中心商 $Q=A/(z_1,\ldots,z_4)$ のこれらの次数での次元は
\[
             (q_0,\ldots,q_7)=(1,4,10,20,31,40,44,40)
\]
である。各 $Q_{p,n}$ の任意の基底を経路に持ち上げると、これらに中心単項式を
掛けたものが七次まで各 $A_{p,n}$ の基底となる。特に六次は84次元、七次は120次元である。
\end{theorem}
\begin{proof}
すでに四中心上の自由性が全次数で知られた窓には結論がある。残りは
第\ref{sec:rank-fifteen-quadratic-parabolic-pencils}節の系より、3上の付値で
$X\ne0,1$, $Y\ne0,(1-X)^2$ の単元域に入る。代数閉体上で始める。
前節の四中心を持つ剰余表示を取り、第\ref{sec:rank-fifteen-integral-annulus}節の
モニックな座標方程式でその表示を持ち上げる。五次までの未切断経路は、
一般点の下界と剰余での同じ次元によって単元主元の語基底を持つ。
この基底で四次中心の方程式を作ると剰余で核は五次元であり、
命題\ref{prop:fixed-residual-regularity}で用いた正則一般窓の四次中心から来る核次元の下界も五である。従って主元小行列式が単元のまま
自由列を選び、指定された四つの中心を同時に持ち上げられる。
各有限長さの中心商は語と関係の有限行列の余核だから、分数体で
$\dim Q_{p,n}\le q_n$ ($0\le n\le7$) となる。

形式的な次数4の変数を持つ多項式環を $S=k[T_1,\ldots,T_4]$ とし、$T_i$ を
$z_i$ として作用させる。$S_{4j}$ の次元は $s_j=\binom{j+3}{3}$ である。
任意の中心商の基底の持上げは各始点の経路を $S$ 上で生成する。実際、商における
像を差し引いた後は中心倍の和であり、残る経路の長さは4だけ短い。中心は次数4で
周期の始点に戻り、右に掛けるとこの生成を始点ごとに使える。
定理\ref{thm:rank-fifteen-generic-levels}の正則一般窓への比較による未切断経路の
下界は、各通常窓で $\dim A_{p,n}\ge h_n$ である。一方、
\[
 h_n\le \dim A_{p,n}\le\sum_{0\le4j\le n}s_j\dim Q_{p,n-4j}
 \le\sum_{0\le4j\le n}s_j q_{n-4j}=h_n
 \qquad(0\le n\le7).
\]
最後は $(1+t+t^2+t^3)^4/(1-t^4)^4=(1-t)^{-4}$ の係数比較である。
全て等号なので商の次元も確定し、生成する単項式は線型独立である。
体拡大を使った場合、元の体上の中心空間は核の体拡大と可換である。
四式の選択空間で七次までの商次元の上界は有限個の開条件をなし、非空である。
標数零の体は無限なので元の体でも同時に選べる。
\end{proof}
\begin{corollary}
この四中心を固定すれば、八次から十一次の商には四始点のそれぞれで下界
$31,20,10,4$ がある。これらの次数のいずれかで対応する上界が得られれば、
その次数で商次元、未切断経路次元、中心単項式基底も同時に確定する。
\end{corollary}
\begin{proof}
$n=8,9,10,11$ なら $n-4,n-8$ はすでに定理の範囲にある。
上の生成による不等式から
$\dim Q_{p,n}\ge h_n-4q_{n-4}-10q_{n-8}$ となる。右辺が掲げた四数である。
商次元がこの下界に達すれば、再び生成写像の両辺の次元が等しくなる。
\end{proof}
